\section*{Capítulo \ref{ch:b2:frontier-orbitals} --- Orbitales frontera y reactividad}

\begin{solution}{exo:b2:frontier-orbitals:1}
Eteno: \omterm{def:b2:frontier-orbitals:frontier}{HOMO} $\alpha + \beta$, \omterm{def:b2:frontier-orbitals:frontier}{LUMO} $\alpha - \beta$. Anión alilo (cuatro electrones): \omterm{def:b2:frontier-orbitals:frontier}{HOMO} $\alpha$ (el
nivel no enlazante), \omterm{def:b2:frontier-orbitals:frontier}{LUMO} $\alpha - 1.414\beta$. Butadieno: \omterm{def:b2:frontier-orbitals:frontier}{HOMO} $\alpha + 0.618\beta$, \omterm{def:b2:frontier-orbitals:frontier}{LUMO} $\alpha -
0.618\beta$.
\end{solution}

\begin{solution}{exo:b2:frontier-orbitals:2}
Duros: \ce{OH-}, \ce{H+}, \ce{F-}. Blandos: \ce{I-}, \ce{RS-}, el carbono del yodometano.
\end{solution}

\begin{solution}{exo:b2:frontier-orbitals:3}
Reaccionan el butadieno (confórmero s-cis accesible) y el ciclopentadieno (fijado en s-cis);
el penta-1,4-dieno no está conjugado y no reacciona. El \omterm{def:b2:frontier-orbitals:diels-alder}{dieno} (2\textit{Z},4\textit{Z}) solo puede
alcanzar la forma s-cis empujando sus dos grupos metilo internos el uno contra el otro:
reacciona muy mal.
\end{solution}

\begin{solution}{exo:b2:frontier-orbitals:4}
Ciclohex-3-eno-1-carbonitrilo: un anillo de seis eslabones, con el doble enlace entre los antiguos C2
y C3 del \omterm{def:b2:frontier-orbitals:diels-alder}{dieno} y el grupo \ce{C#N} sobre uno de los dos carbonos del anillo procedentes del
\omterm{def:b2:frontier-orbitals:diels-alder}{dienófilo}.
\end{solution}

\begin{solution}{exo:b2:frontier-orbitals:5}
$\Delta\varepsilon = \qty{8}{eV}$: perturbativa, $2\beta^2/\Delta\varepsilon = \qty{0.25}{eV}$. Exacta: nivel inferior
$-6 - \sqrt{16 + 1} = \qty{-10.123}{eV}$, ganancia $2 \times 0.123 = \qty{0.246}{eV}$.
\end{solution}

\begin{solution}{exo:b2:frontier-orbitals:6}
El yodometano es un \omterm{def:b2:frontier-orbitals:hard-soft}{electrófilo blando}: \omterm{def:b2:frontier-orbitals:control}{control orbital}, ataque por el átomo con el mayor
coeficiente en el \omterm{def:b2:frontier-orbitals:frontier}{HOMO}, el carbono. Un carbocatión libre es un \omterm{def:b2:frontier-orbitals:hard-soft}{electrófilo duro}: \omterm{def:b2:frontier-orbitals:control}{control de carga},
ataque en parte por el nitrógeno, que lleva más carga negativa.
\end{solution}

\begin{solution}{exo:b2:frontier-orbitals:7}
El C4 del \omterm{def:b2:frontier-orbitals:diels-alder}{dieno} (mayor coeficiente en el \omterm{def:b2:frontier-orbitals:frontier}{HOMO}) se enlaza al carbono $\beta$ del propenal (mayor
coeficiente en el \omterm{def:b2:frontier-orbitals:frontier}{LUMO}): 2-metoxiciclohex-3-eno-1-carbaldehído, con los dos sustituyentes en
carbonos vecinos (producto ``orto'').
\end{solution}

\begin{solution}{exo:b2:frontier-orbitals:8}
Un esqueleto de biciclo[2.2.1]hepteno (norborneno) con el grupo éster sobre un carbono del
antiguo \omterm{def:b2:frontier-orbitals:diels-alder}{dienófilo}, orientado hacia el nuevo enlace \ce{C=C}, bajo el anillo de seis eslabones
(endo), del lado opuesto al puente \ce{CH2}.
\end{solution}

\begin{solution}{exo:b2:frontier-orbitals:9}
Los dos grupos carbonilo rebajan el \omterm{def:b2:frontier-orbitals:frontier}{LUMO} del \omterm{def:b2:frontier-orbitals:diels-alder}{dienófilo}: su diferencia con el \omterm{def:b2:frontier-orbitals:frontier}{HOMO} del
ciclopentadieno es mucho menor que la del eteno, la estabilización $\beta^2/\Delta\varepsilon$ del
estado de transición es mayor y la barrera, más baja.
\end{solution}

\begin{solution}{exo:b2:frontier-orbitals:10}
La reacción es estereoespecífica: el maleato da el aducto con los dos grupos éster cis
(endo,endo como producto mayoritario; es aquiral, un compuesto meso); el fumarato da el
aducto trans, con un éster endo y otro exo, formado como una pareja de enantiómeros (racémico).
\end{solution}

\begin{solution}{exo:b2:frontier-orbitals:11}
Dos enlaces $\pi$ se convierten en dos enlaces $\sigma$ más fuertes: $\Delta_r H^\circ < 0$. Dos moléculas se convierten en una:
$\Delta_r S^\circ < 0$. $\Delta_r G^\circ = \Delta_r H^\circ - T\Delta_r S^\circ$ crece con $T$ y se vuelve
positiva por encima de $T_i = \Delta_r H^\circ/\Delta_r S^\circ$: a alta temperatura el aducto se rompe de nuevo
(retro-Diels--Alder).
\end{solution}

\begin{solution}{exo:b2:frontier-orbitals:12}
Dos orbitales llenos de pares no enlazantes, uno junto a otro, forman una interacción a cuatro electrones, que es
repulsiva. La molécula gira alrededor del enlace \ce{N-N} hasta que los pares no enlazantes son aproximadamente
perpendiculares (una conformación gauche), donde su solapamiento, y la repulsión, son
pequeños.
\end{solution}

\begin{solution}{pb:b2:frontier-orbitals:1}
\textbf{1.} Un anillo de cinco eslabones con dos dobles enlaces conjugados y un \ce{CH2}: el anillo
mantiene la unidad de \omterm{def:b2:frontier-orbitals:diels-alder}{dieno} en la conformación s-cis.
\textbf{2.} El \omterm{def:b2:frontier-orbitals:diels-alder}{dienófilo}, a través de uno de sus dobles enlaces.
\textbf{3.} Un esqueleto de norborneno fusionado con un anillo de ciclopenteno; los dos nuevos enlaces $\sigma$ unen
los extremos del \omterm{def:b2:frontier-orbitals:diels-alder}{dieno} con los dos carbonos del doble enlace del \omterm{def:b2:frontier-orbitals:diels-alder}{dienófilo}.
\textbf{4.} El \omterm{def:b2:frontier-orbitals:endo}{aducto endo} (\omterm{def:b2:frontier-orbitals:endo}{regla endo}): el anillo restante del \omterm{def:b2:frontier-orbitals:diels-alder}{dienófilo} queda bajo
el \omterm{def:b2:frontier-orbitals:diels-alder}{dieno} en el estado de transición.
\textbf{5.} Sí: es concertada, ambos enlaces se forman por la misma cara de cada participante y
se conserva la geometría del \omterm{def:b2:frontier-orbitals:diels-alder}{dienófilo}.
\textbf{6.} El ciclopentadieno es a la vez un buen \omterm{def:b2:frontier-orbitals:diels-alder}{dieno} (fijado en s-cis) y un \omterm{def:b2:frontier-orbitals:diels-alder}{dienófilo}, y la
reacción es concertada, sin ningún intermedio cargado que estabilizar.
\textbf{7.} Negativo: dos enlaces $\pi$ se sustituyen por dos enlaces $\sigma$.
\textbf{8.} Negativo: dos moléculas dan una.
\textbf{9.} $\Delta_r G^\circ = \Delta_r H^\circ - T\Delta_r S^\circ$ con ambos términos negativos: negativa a
$T$ baja, nula en $T_i = \Delta_r H^\circ/\Delta_r S^\circ$ y positiva por encima.
\textbf{10.} A temperatura ambiente, $T < T_i$: la dimerización está favorecida, aunque es lenta. Cerca
del punto de ebullición del dímero el equilibrio favorece al monómero, y se alcanza
deprisa.
\textbf{11.} El monómero sale de la mezcla caliente como vapor: retirar un producto mantiene
el equilibrio desplazándose hacia el monómero (Le Chatelier).
\textbf{12.} A temperatura ambiente se dimeriza de nuevo; el frío frena esa reacción.
\textbf{13.} $0.62 - (-0.62) = 1.24|\beta|$.
\textbf{14.} $0.62 - (-0.35) = 0.97|\beta|$.
\textbf{15.} El ciclopentadieno con el anhídrido maleico: la menor diferencia de energía da una mayor
estabilización del estado de transición, y una reacción más rápida.
\textbf{16.} Están conjugados con el enlace \ce{C=C} y son atractores de electrones: como en el
propenal, hacen bajar el nivel $\pi^*$.
\textbf{17.} Su \omterm{def:b2:frontier-orbitals:frontier}{LUMO}, que tiene lóbulos sobre los carbonos carbonílicos enfrentados a los carbonos centrales del
\omterm{def:b2:frontier-orbitals:frontier}{HOMO} del \omterm{def:b2:frontier-orbitals:diels-alder}{dieno}.
\textbf{18.} Esqueleto de norborneno; el anillo de anhídrido bajo el nuevo doble enlace, del lado opuesto al
puente \ce{CH2}.
\textbf{19.} Los dos carbonos cabeza de puente y los dos carbonos que llevan el anhídrido: cuatro
estereocentros.
\textbf{20.} No: un plano de simetría pasa por el puente \ce{CH2} y entre las dos mitades
de la molécula; es un compuesto meso.
\textbf{21.} Eran cis en el \omterm{def:b2:frontier-orbitals:diels-alder}{dienófilo}, y la adición concertada los mantiene cis.
\textbf{22.} El mismo esqueleto con el anillo de anhídrido del lado del puente \ce{CH2}.
\textbf{23.} No: pasar de endo a exo supondría romper y volver a formar enlaces \ce{C-C}, mediante
la reacción inversa, que exige calentar.
\textbf{24.} El agua abre el anhídrido: el ácido dicarboxílico cis, con los dos grupos \ce{COOH} en endo.
\textbf{25.} El \omterm{def:b2:frontier-orbitals:endo}{aducto endo} tiene $\boldsymbol{4}$ estereocentros.
\end{solution}
